\section{假设检验}

本节简单介绍假设检验的核心思想和相关概念。

本节要点：
\begin{itemize}
    \item 了解假设检验。
\end{itemize}

%============================================================
\subsection{假设的提法}

\begin{definition}[假设的提法]
我们称检验问题中两个对立的命题为{\bf 假设}或{\bf 统计假设}，并将其中一个命题称为{\bf 原假设}或{\bf 零假设}，记作$H_0$，另一个命题称为{\bf 备选假设}或{\bf 对立假设}，记作$H_1$，于是检验问题可记作$\left( H_0,H_1 \right) $。
\end{definition}

例1的假设可写成：
\[
H_0:\mu =500, \quad H_1:\mu \ne 500
\]

例2的假设可写成：
\[
H_0:\text{硬币没有问题}, \quad H_1:\text{硬币不均匀}
\]

%============================================================
\subsection{检验的基本原理}

以步骤说明检验的基本原理：
\begin{enumerate}
    \item 假设总体的特征，描述检验问题$\left( H_0,H_1 \right) $；
    \item 抽样得到样本$\left( X_1,\cdots ,X_n \right) $；
    \item 选取合适的统计量$G$进行计算，获得观察值$g=g\left( x_1,\cdots ,x_n \right) $；
    \item 根据该统计量的分布和计算得到的观察值，计算该值出现的概率$P$；
    \item 根据概率判断本次抽样是大概率事件还是小概率事件；
    \item 若是大概率事件，我们接受对总体的假设$H_0$，反之，我们接受$H_1$。
\end{enumerate}

不难发现，假设检验的核心有两点：
\begin{itemize}
    \item 合适的统计量$G$，$G$必须能正确描述样本和总体的差异的随机性；
    \item 事件发生的概率大小的阈值，该阈值应该能正确反映差异是由样本的随机性造成，还是样本和总体的偏差造成。
\end{itemize}

同时也不难发现，假设检验的问题在于一切皆概率，小概率事件还是会发生的，无论接受$H_0$还是$H_1$，都有可能是错的。

%============================================================
\subsection{假设检验的两类错误}

\begin{definition}[拒绝域和接受域]
假设检验中，若小概率事件A发生，则拒绝$H_0$，我们称导致小概率事件A发生的统计量$g\left( X_1,\cdots ,X_n \right) $的全体取值范围为{\bf $H_0$的拒绝域}，常记作$W$。
同样，若小概率事件A不发生，则接受$H_0$，我们称导致小概率事件A不发生的统计量$g\left( X_1,\cdots ,X_n \right) $的全体取值范围为{\bf $H_0$的接受域}。
\end{definition}

\begin{definition}[两类错误]
当真实情况$H_0$成立，但检验结果拒绝$H_0$，我们称之为{\bf 第一类错误}或{\bf 弃真错误}，发生概率称为{\bf 显著性水平}，常记作$\alpha $。
反之，当真实情况是$H_1$成立，但检验结果拒绝$H_1$，我们称之为{\bf 第二类错误}或{\bf 存伪错误}，发生概率常记作$\beta $。
\end{definition}

显著性水平$\alpha $代表了样本和总体的差异从抽样的随机性到怀疑假设这一质变，表现为统计量的密度函数上的上侧分位点。
实际中我们一般取$\alpha =0.05$，根据统计量$G$的具体表达式查到上侧分位点$g_{\alpha =0.025}$（若密度函数为偶函数），确定拒绝域$W=\left( -g_{\alpha =0.025},g_{\alpha =0.025} \right) $，当统计观察值$g\notin W$，我们接受$H_0$。

\begin{table}[h]
\centering
\begin{tabular}{ccc}
    \toprule
    & 检验结果接受$H_0$ & 检验结果接受$H_1$\\
    \midrule
    真实情况$H_0$成立 & 判断正确 & 第一类错误（弃真错误）\\
    真实情况$H_1$成立 &  第二类错误（存伪错误） & 判断正确\\
    \bottomrule
\end{tabular}
\end{table}

一般来讲，当我们固定样本容量时$\alpha ,\beta $是一对矛盾，而第一类错误的危害普遍认为要比第二类错误大，所以策略上我们通常采取“先限制$\alpha $再降低$\beta $”。
理论证明，选用合适的统计量可有效降低$\beta $。

%============================================================
\subsection{假设检验的步骤}

我们可总结假设检验的步骤：
\begin{enumerate}
    \item 建立检验问题$\left( H_0,H_1 \right) $；
    \item 根据问题选择合适的统计量$G$，并确定其分布，该步为重点；
    \item 根据显著性水平$\alpha $（一般取0.05，也可根据实际调整）和$G$的分布确定$H_0$的拒绝域$W$，该步为难点；
    \item 根据抽到的样本值$\left( x_1,\cdots ,x_n \right) $计算观察值$g=g\left( x_1,\cdots ,x_n \right) $；
    \item 若$g\notin W$，我们接受$H_0$，反之若$g\in W$，我们拒绝$H_0$。
\end{enumerate}




